\documentclass[11pt,a4paper]{amsart}

\usepackage{amsmath,amssymb,amsthm}
\usepackage{mathtools}
\usepackage{booktabs}
\usepackage{array}
\usepackage{hyperref}
\usepackage{xcolor}
\usepackage[margin=1in]{geometry}

\hypersetup{
  colorlinks=true,
  linkcolor=blue!50!black,
  urlcolor=blue!60!black,
  citecolor=blue!50!black
}

\theoremstyle{plain}
\newtheorem{theorem}{Theorem}
\newtheorem{proposition}[theorem]{Proposition}
\newtheorem{corollary}[theorem]{Corollary}
\newtheorem{lemma}[theorem]{Lemma}

\theoremstyle{definition}
\newtheorem{definition}[theorem]{Definition}
\newtheorem{remark}[theorem]{Remark}
\newtheorem{observation}[theorem]{Observation}

\DeclareMathOperator{\Aut}{Aut}
\DeclareMathOperator{\Vec}{Vec}
\DeclareMathOperator{\Rep}{Rep}
\DeclareMathOperator{\Mod}{Mod}
\DeclareMathOperator{\Irr}{Irr}
\DeclareMathOperator{\Hom}{Hom}
\DeclareMathOperator{\Ext}{Ext}
\DeclareMathOperator{\Tor}{Tor}
\DeclareMathOperator{\im}{im}
\DeclareMathOperator{\Inn}{Inn}
\DeclareMathOperator{\Out}{Out}

\newcommand{\GrpCenter}[1]{Z\!\left(#1\right)}
\newcommand{\Drinfeld}[2]{Z\!\left(\Vec_{#1}^{#2}\right)}
\newcommand{\trivcoc}{\mathbf{1}}
\newcommand{\ord}{\mathrm{ord}}
\newcommand{\hcl}[1]{[#1]}

\title{The trivial and order--$2$ twists of $Q_8$\\
       give inequivalent modular tensor categories\\
       but identical fusion rings}
\author{}
\date{\today}

\begin{document}
\maketitle

\begin{abstract}
We compare the modular data of the two twisted Drinfeld doubles
$\Drinfeld{Q_8}{\trivcoc}$ and $\Drinfeld{Q_8}{\omega_4}$, where
$\omega_4$ is the unique order--$2$ class in
$H^3(Q_8;U(1))\cong\mathbb Z/8$. By an exhaustive permutation search we
show that there is no relabelling of simple objects identifying the two
modular data $(S_i,T_i)$, even after a Galois twist $i\leftrightarrow -i$,
so $\Drinfeld{Q_8}{\trivcoc}$ and $\Drinfeld{Q_8}{\omega_4}$ are
inequivalent as modular tensor categories. Combined with
Mignard--Schauenburg's theorem that modular data is a complete invariant
for $|G|<32$, this implies $\Vec_{Q_8}$ and $\Vec_{Q_8}^{\omega_4}$ are
not categorically Morita equivalent. \emph{Nevertheless}, their fusion
rings (Grothendieck rings) coincide: we exhibit an explicit ring
isomorphism. The entire numerical difference is concentrated in a single
$g$--twisted sector, where two $T$--eigenvalues swap. Hence Evans--Gannon's
``extension type'' property cannot be detected from the fusion ring
alone, but is visible from the joint $(S,T)$ data.
\end{abstract}

\section{Setup and question}

For a finite group $G$ and a $3$--cocycle
$\omega\in Z^3(G;U(1))$ the twisted Drinfeld double
$\Drinfeld{G}{\omega}=Z(\Vec_G^\omega)$ is a modular tensor category
(MTC) whose simple objects are pairs $(g,\rho)$ with $g$ ranging over
conjugacy class representatives of $G$ and $\rho$ a projective
irreducible representation of the centraliser $C_G(g)$ with respect to a
$2$--cocycle $\beta_g$ obtained from $\omega$
by the formula
\begin{equation}
\beta_g(h,k) = \omega(g,h,k) + \omega(h,k,g) - \omega(h,g,k).
\label{eq:beta}
\end{equation}

Evans--Gannon \cite{EvansGannon} call a (type--$1$, equivalently:
extension--type) module category over $\Drinfeld{G}{\omega}$
\emph{trivialising} when its category of local modules is
$\Vec_{\mathbb C}$. Their Corollary~2 classifies such trivialising
extensions by pairs $(K\le G,\;\psi'\in Z^2(K;U(1)))$ with
$\omega|_K$ a coboundary. For $\omega=\trivcoc$ this is automatic with
$K=G$; for nontrivial $\omega$ one needs a proper subgroup $K$ on which
$\omega$ trivialises, and the supply of such $K$ is what
Evans--Gannon call ``the obstruction''.

It is folklore in the physics literature that this obstruction is
detectable from modular data $(S,T)$ alone. The user--posed question:

\begin{quote}
  \emph{Does $Q_8$ provide a counterexample? That is, do the trivial
  cocycle and the unique order--$2$ cocycle on $Q_8$ produce
  modular data that agree under some natural equivalence relation?}
\end{quote}

\section{Verified setup}

We use Gruen's GAP implementation (\cite{Gruen}, in
\verb|DRINFELD/modular_data.g|) of the Mignard--Schauenburg /
Gruen--Morrison algorithms for computing $(S,T)$, applied via the
\textsf{HAP} package.

\begin{itemize}
\item \verb|SmallGroup(8,4)| has GAP \verb|StructureDescription "Q8"|;
      its element--order distribution is $\{1\colon 1,\ 2\colon 1,\ 4\colon 6\}$,
      so the group is the quaternion group $Q_8$.
\item Gruen's universal--coefficients computation returns
      \verb|UCT.exponent = 8| and \verb|UCT.torsion = [8]|, recovering
      $H^3(Q_8;U(1))\cong\mathbb Z/8$.
\item $\Aut(Q_8)$ acts trivially on $H^3(Q_8)$: there are exactly $8$
      orbits, each a singleton. The eight orbit representatives stored in
      \verb|modular_data/8/4/3-cocycles.txt| have orders
      $\{1, 2, 4, 4, 8, 8, 8, 8\}$ in $\mathbb Z/8$, accounting for all
      $8$ elements.
\item File index $1$ has cocycle vector $[0,0,0]$ whose
      $H^3$--class is $[0,0]$ --- the \textbf{trivial cocycle} $\trivcoc$.
\item File index $2$ has cocycle vector $[0,4,0]$ whose $H^3$--class
      is $[0,4]$, the unique class of order $2$ --- write this
      $\omega_4$.
\item Both modular data sets have rank $22$ ($=5+5+4+4+4$, from the
      conjugacy class / centraliser--irrep counts of $Q_8$).
\item A fresh \texttt{ComputeModularData} run reproduces each stored
      $(S_i,T_i)$ up to permutation; the random row selection in
      \texttt{GenerateModularData} only reorders simples.
\end{itemize}

Throughout, $(S_1,T_1)$ denotes the modular data stored as cocycle~$1$
(trivial) and $(S_2,T_2)$ the modular data stored as cocycle~$2$ (order
$2$).

\section{Equivalence ladder}

We tested whether $(S_1,T_1)$ and $(S_2,T_2)$ are identified by the
following equivalence relations, in increasing order of strength.

\begin{enumerate}
\item \emph{$T$--spectrum equality} (multiset of diagonal entries).
\item \emph{Per--row $S$ multiset equality}.
\item \emph{Fusion ring isomorphism} (Grothendieck rings agree up to a
       relabelling of simples).
\item \emph{Joint (Fusion ring, $T$--vector) equivalence}: a single
       permutation aligning both $N$ and $T$.
\item \emph{Modular data equivalence}: a single permutation aligning
       both $S$ and $T$.
\item \emph{Galois--twisted modular data equivalence}: the same plus an
       overall complex--conjugation symmetry $i\leftrightarrow -i$.
\end{enumerate}

The $S$--matrices entries lie in $\mathbb Q$ in this case, so the Galois
action is trivial on $S$ and acts only on $T$ (which has entries in
$\mathbb Q(i)$).

\medskip
\noindent
The fusion coefficients are computed from $S$ via the Verlinde formula
\begin{equation}
N_{ij}^k = \sum_{a} \frac{S_{ia}\,S_{ja}\,\overline{S_{ka}}}{S_{0a}}.
\label{eq:verlinde}
\end{equation}

We computed $N_1$ and $N_2$ from $S_1, S_2$ and verified their entries
are nonnegative integers, as expected. For each candidate equivalence we
ran an exact backtracking search across the $22\times 22$ index space,
pruning by the natural per--row signature.

\begin{table}[h!]
\centering
\renewcommand{\arraystretch}{1.15}
\begin{tabular}{lc}
\toprule
Equivalence relation & $(S_1,T_1) \sim (S_2,T_2)?$ \\
\midrule
$T$--spectrum (multiset) & \textbf{Yes} \\
$S$--row multisets & \textbf{Yes} \\
Fusion ring (Grothendieck ring) & \textbf{Yes} \\
Joint (fusion ring,\,$T$--vector) up to permutation & \textbf{No} \\
Full $(S,T)$ up to permutation & \textbf{No} \\
$(S,T)$ up to permutation $+$ Galois twist & \textbf{No} \\
\bottomrule
\end{tabular}
\caption{Equivalence ladder for $(S_1,T_1)$ vs.\ $(S_2,T_2)$.}
\end{table}

\paragraph*{Common $T$--spectrum.}
Both $T_1$ and $T_2$ have multiset of diagonal entries
\[
  \bigl\{1^{\times 12},\ (-1)^{\times 4},\ i^{\times 3},\ (-i)^{\times 3}\bigr\}.
\]

\paragraph*{Fusion ring isomorphism.}
The backtracking search returns the explicit permutation
\begin{equation}
\pi_F \;=\; (6,8)(7,9)\,(10,12)(11,13)\,(15,16,17)\,(19,21)(20,22),
\label{eq:piF}
\end{equation}
of order $6$, satisfying
$N_2[i,j,k] = N_1[\pi_F(i),\pi_F(j),\pi_F(k)]$ for all $i,j,k$. So
the Grothendieck rings of $\Drinfeld{Q_8}{\trivcoc}$ and
$\Drinfeld{Q_8}{\omega_4}$ are isomorphic.

\paragraph*{Modular data inequivalence.}
The exact backtracking search for permutations $\pi$ satisfying
$T_2[i] = T_1[\pi(i)]$ and $S_2[i,j]=S_1[\pi(i),\pi(j)]$ exhausts the
search space and returns no solution. The same search with a Galois
twist $T_2 \mapsto \overline{T_2}$ (and the corresponding twist on $S$,
which acts trivially since $S_2\in\mathbb Q^{22\times 22}$) also returns
no solution. Sanity--checks confirm the search recovers the identity
permutation when comparing $(S_1,T_1)$ to itself and the correct
permutation when comparing to a known random permutation of itself.

\paragraph*{Joint $(N,T)$ inequivalence.}
We also enumerated permutations $\pi$ satisfying \emph{both}
$N_2[i,j,k]=N_1[\pi(i),\pi(j),\pi(k)]$ and $T_2[i]=T_1[\pi(i)]$. The
count is $0$. As a clean discriminating invariant,
\[
  |\Aut(N_1, T_1)| = 6,\qquad |\Aut(N_2, T_2)| = 24,
\]
where $\Aut(N,T)$ denotes permutations of $\{1,\dots,22\}$ preserving
both $N$ and $T$. Since the two automorphism counts differ, the labelled
pairs $(N_1,T_1)$ and $(N_2,T_2)$ are not isomorphic, even though the
underlying fusion rings are.

\paragraph*{Note on the heuristic.}
Gruen's \texttt{BasicEquivalentModularDataTests} returns \texttt{true}
on this pair. The heuristic only checks per--$T$--block sorted
$S$--diagonal multisets and does not verify cross--row consistency, so
it is genuinely fooled here. The exact search above settles the
question.

\section{Where exactly do the two modular data diverge?}

We label each row of the modular data with its
$(\text{conjugacy class},\;\dim\rho)$ pair by re--running
\texttt{GenerateSimpleObjects} (deterministic) and aligning to the
stored ordering with the same backtracking search.

The simples partition by conjugacy class as follows. Write
$c_{(\trivcoc)}=\{1\}$, $c_{(-1)}=\{-1\}$, and $c_A,c_B,c_C$ for the
three order--$4$ conjugacy classes (each of size $2$).

\begin{table}[h!]
\centering
\renewcommand{\arraystretch}{1.2}
\begin{tabular}{l c c c c}
\toprule
class & rows in stored data & $\dim$ & $T$ values (data~$1$) & $T$ values (data~$2$) \\
\midrule
$c_{(\trivcoc)}$ & $1,2,3,4,5$    & $1,1,1,1,2$ & $1,1,1,1,1$       & $1,1,1,1,1$ \\
$c_{(-1)}$       & $14,15,16,17,18$ & $1,1,1,1,2$ & $1,1,1,1,-1$       & $1,1,1,1,-1$ \\
$c_A$            & $6,7,8,9$      & $1,1,1,1$   & $1,-1,i,-i$       & $1,-1,i,-i$ \\
$c_B$            & $10,11,12,13$  & $1,1,1,1$   & $1,-1,i,-i$       & $1,-1,i,-i$ \\
$c_C$            & $19,20,21,22$  & $1,1,1,1$   & $\boldsymbol{1,-1},i,-i$ & $\boldsymbol{-1,1},i,-i$ \\
\bottomrule
\end{tabular}
\caption{Sector decomposition of the simples and their twists.}
\label{tab:sectors}
\end{table}

\begin{observation}
\label{obs:onesector}
The entire numerical difference between the two modular data is a swap
of two $T$--eigenvalues in exactly one $g$--twisted sector --- the third
order--$4$ conjugacy class $c_C$. The other four sectors
($c_{(\trivcoc)},c_{(-1)},c_A,c_B$) are identical between the two cocycles.
At the level of the diff between the stored files
\verb|1.txt| and \verb|2.txt|, this corresponds to a transposition of
$T$--diagonal positions $19$ and $20$, accompanied by all the necessary
sign changes in the rows / columns of $S$ that link sector $c_C$ to the
other sectors.
\end{observation}

\section{Interpretation}

\subsection*{Why fusion rings agree but $T$ differs.}
Each cyclic centraliser $C_G(g)\cong\mathbb Z/4$ for $g\in c_A,c_B,c_C$
has trivial Schur multiplier, so the $2$--cocycle $\beta_g$ from
\eqref{eq:beta} is necessarily a coboundary. Hence the projective
characters appearing in the simple--object enumeration are still
ordinary characters of $\mathbb Z/4$: dimensions are unaffected.
\emph{Fusion rules see only dimensions and multiplicities and are
phase--blind}, so the fusion ring is unchanged by the cocycle.

The explicit choice of $1$--cochain trivialising $\beta_g$ does depend
on $g$, however, and rotates the $T$--eigenvalues
$\rho(g)/\dim\rho$ by a $g$--dependent phase. For $\omega=\omega_4$ on
$Q_8$ this phase happens to be trivial on two of the three order--$4$
classes and equal to $-1$ on the third (Observation
\ref{obs:onesector}). The resulting $T$--diagonal therefore differs from
the trivial--cocycle one only in sector $c_C$, by a swap of two entries.

\subsection*{Implication for ``extension type''.}
By Mignard--Schauenburg \cite[Theorem on modular data
completeness]{MignardSchauenburg}, modular data is a complete invariant
of $\Drinfeld{G}{\omega}$ for $|G|<32$. Combined with our explicit
non--existence of a permutation aligning $(S_1,T_1)$ with $(S_2,T_2)$,
we conclude:

\begin{proposition}
$\Drinfeld{Q_8}{\trivcoc}$ and $\Drinfeld{Q_8}{\omega_4}$ are
inequivalent as modular tensor categories. Equivalently
$\Vec_{Q_8}$ and $\Vec_{Q_8}^{\omega_4}$ are not categorically Morita
equivalent.
\end{proposition}

In particular, $Q_8$ is \textbf{not} a counterexample to the folklore
``modular data detects $[\omega]=0$'' under the standard reading of
``modular data'' as the pair $(S,T)$.

\medskip
\noindent
However, the analysis surfaces a clean weaker statement.

\begin{observation}
The fusion ring and the $T$--multiset, taken \emph{individually}, do not
distinguish $\Drinfeld{Q_8}{\trivcoc}$ from $\Drinfeld{Q_8}{\omega_4}$:
the fusion rings are isomorphic via $\pi_F$, and the $T$--multisets are
equal. The discriminator that detects $[\omega_4]\ne 0$ is the joint
$(N,T)$ structure, equivalently any invariant that ties a $T$--eigenvalue
to its position in the fusion algebra (such as the full $S$).
Anyone reading ``extension type'' off the fusion ring alone --- or off
the $T$--spectrum alone --- would conflate the trivial and order--$2$
twists of $Q_8$.
\end{observation}

\section*{Reproducibility}

The investigation is reproducible from the
\href{https://github.com/thomas-gemunden/DRINFELD}{DRINFELD repository}.
Required software: GAP $\ge 4.14$ with the \textsf{HAP} and \textsf{IO}
packages. A standalone GAP $4.15.1$ install (via the
\verb|gap-system/gap| Homebrew tap) and Sage's bundled GAP both work.

The GAP scripts driving the searches live alongside this file in
\verb|scripts/|:

\begin{itemize}\itemsep0pt
\item \verb|verify_q8.g| --- confirms \verb|SmallGroup(8,4)=Q_8| and
      identifies the two cocycles by $H^3$ class and order.
\item \verb|recompute_check.g|, \verb|recompute_perm.g| --- recompute
      $(S,T)$ from scratch and verify the stored data is correct up to
      permutation.
\item \verb|find_permutation.g|, \verb|sanity_check.g| --- exact
      backtracking search for permutations aligning $(S,T)$.
\item \verb|galois_and_fusion.g| --- Galois--twisted search; Verlinde
      fusion ring; per--row signature comparison.
\item \verb|fusion_plus_T.g|, \verb|sanity_count.g| --- joint
      $(N,T)$ search and the $6$--vs--$24$ automorphism--count
      discriminator.
\item \verb|label_simples.g| --- recovers $(g,\rho)$ labels for every
      row of stored \verb|1.txt|/\verb|2.txt| and prints $\pi_F$ in
      label form.
\end{itemize}

\noindent
To run them, \verb|cd| into \verb|scripts/| and invoke
\verb|./setup.sh| once (this creates the \verb|Modular_Data| symlink
that Gruen's \verb|LoadModularData| expects), then
\verb|gap -q -b -c 'Read("<script>.g");'|. The actual modular data
files live outside this repository (the full database is $\sim 77$\,GB)
at the sibling path \verb|../../../modular_data/8/4/{1,2}.txt|.

\begin{thebibliography}{9}

\bibitem{EvansGannon}
D.~E.~Evans and T.~Gannon,
\emph{Reconstruction and local extensions for twisted group doubles,
and permutation orbifolds},
arXiv:\href{https://arxiv.org/abs/1804.11145}{1804.11145}.

\bibitem{Gruen}
A.~Gruen,
\emph{Computing Modular Data for Drinfeld Centres of Pointed Fusion
Categories},
B.Phil.\ thesis, Australian National University, 2017.
Available at
\href{https://tqft.net/web/research/students/AngusGruen/thesis.pdf}%
{tqft.net/web/research/students/AngusGruen/thesis.pdf}.

\bibitem{MignardSchauenburg}
M.~Mignard and P.~Schauenburg,
\emph{Morita equivalence of pointed fusion categories of small rank},
arXiv:\href{https://arxiv.org/abs/1708.06538}{1708.06538}.

\end{thebibliography}

\end{document}
